% chapters/theory/02.tex

\chapter{Modern Transformer Design and Model Families}
\label{chap:modern-transformers-and-families}

The original Transformer architecture is remarkably stable: attention blocks, feedforward blocks, residual connections, and normalization remain the backbone of most language models. At the same time, many details have evolved in response to practical constraints such as long context, training stability, and inference efficiency. In parallel, three major model families emerged—encoder--decoder, encoder-only, and decoder-only—each shaped by different objectives and use cases. This chapter explains the most important modern design choices and then connects them to the best-known model families, with T5 and BERT as anchor points.

\section{What attention heads reveal}
Multi-head attention is often described as “multiple attention computations in parallel,” but it is more useful to think of it as “multiple learned projections of the same sequence.” Each head has its own projection matrices that map token embeddings into queries, keys, and values. Because these projections differ, heads can learn different relationship patterns among tokens. A simple interpretability lens is the attention map: for a token like a pronoun, some heads assign high weight to antecedents (e.g., a nearby noun phrase), reflecting the model’s ability to track reference. Attention maps are not a complete explanation of model behavior, but they provide an intuitive window into what heads may be doing.

A common confusion is whether each head has its own separate MLP. It does not. The head-specific components are the projection matrices for \(Q\), \(K\), and \(V\). The feedforward network is a separate sublayer applied after attention and is shared at the layer level, not duplicated per head.

\section{Positional information: from absolute embeddings to modern methods}
Self-attention alone does not encode token order. If we permute tokens, the attention mechanism has no inherent concept of “before” and “after.” Positional information must be injected, and the design space here has seen some of the most impactful changes since 2017.

\subsection{Learned absolute positional embeddings}
A direct approach assigns an embedding vector to each position \(p \in \{1,\dots,L_{\max}\}\) and adds it to the token embedding. This is conceptually simple and works well when the maximum training length \(L_{\max}\) matches the deployment length. But it has two practical limitations. First, it binds the model to the maximum length seen during training; positions beyond \(L_{\max}\) require extrapolation. Second, it can pick up dataset-specific biases if certain structures repeatedly occur at particular positions.

\subsection{Sinusoidal positional embeddings}
The original Transformer also proposed fixed sinusoidal encodings, where each position is mapped to a vector of sines and cosines at multiple frequencies. The key intuition is that dot products between sinusoidal position vectors become functions of relative distance via trigonometric identities, providing a smooth way to generalize positional relationships beyond the training range. In practice, sinusoidal and learned embeddings can both work, and the choice often depends on training and extrapolation goals.

\subsection{Relative position bias: injecting position into attention logits}
A deeper insight is that for attention, we do not primarily care about absolute position; we care about relative distance between tokens. This motivates injecting positional information directly into the attention score computation. In its simplest form, one adds a bias term \(B(\Delta)\) to the attention logits, where \(\Delta\) is the relative offset between query and key positions. Because softmax normalizes scores, adding a bias is a natural way to encourage the model to attend differently to near vs far tokens.

T5 popularized a learned relative position bias by bucketizing relative distances and learning a bias per bucket. This approach shifts positional reasoning into the attention layer itself, rather than relying on the model to infer it indirectly from positional embeddings added at the input.

\subsection{ALiBi: deterministic linear bias}
ALiBi uses a deterministic distance-dependent bias rather than learned buckets. The guiding idea is to penalize attention to far tokens more than to near tokens through a simple linear form. This can be attractive because it removes learned positional parameters and can extrapolate to longer contexts.

\subsection{Rotary position embeddings (RoPE)}
Modern LLMs frequently use rotary position embeddings. The core idea is to rotate query and key vectors by position-dependent angles before taking dot products. Because the dot product of rotated vectors can be expressed as a function of relative angle differences, RoPE encodes relative position directly into \(QK^\top\). The 2D rotation picture is familiar from linear algebra; in high dimensions, RoPE applies rotations in independent 2D subspaces across the embedding dimension using a frequency schedule. The practical payoff is strong performance and useful extrapolation properties, which is why RoPE (and variants) are widely adopted.

\section{Normalization: where it sits and what form it takes}
Normalization stabilizes training by controlling activation scales. Two changes became common in modern LLMs: where normalization is applied relative to residual blocks, and which normalization variant is used.

\subsection{Post-norm vs pre-norm}
The original Transformer used a “post-norm” structure: add the residual and then normalize. Many modern architectures use “pre-norm”: normalize the input to the sublayer first, then add the residual after the sublayer. Pre-norm tends to improve optimization stability for deep models and is now a common default.

\subsection{LayerNorm vs RMSNorm}
LayerNorm normalizes activations using the mean and variance across dimensions and learns both a scale and a bias. RMSNorm removes the mean-centering step and typically learns only a scale, normalizing by root-mean-square magnitude. In practice RMSNorm often achieves comparable stability with fewer parameters and lower overhead, which is why it appears frequently in modern LLM implementations.

BatchNorm, by contrast, depends on batch statistics and can behave differently at training vs inference; this is one reason it is uncommon in Transformers.

\section{Attention efficiency: sparse attention and sliding windows}
Full self-attention scales as \(O(n^2)\) in sequence length \(n\). This becomes expensive for long contexts. One family of approaches reduces compute by restricting which tokens can attend to which others.

A prominent idea is sliding-window (local) attention, where each token attends only to a neighborhood. Longformer explored variants of this, and later models often mix local and global attention layers. Even when each layer has a limited attention window, stacking layers expands the effective receptive field: information can propagate across the sequence through multiple hops. The analogy to receptive fields in convolutional networks is useful: local operations composed over depth yield broader influence.

Sparse attention is not a universal best choice; it is a trade-off between long-context capability, computational cost, and the model’s ability to combine distant information.

\section{Sharing key/value projections: MHA, GQA, and MQA}
Another efficiency axis concerns how attention heads are parameterized. Standard multi-head attention gives each head its own query, key, and value projections. But during decoding, keys and values are repeatedly reused and stored in the KV cache; their memory footprint becomes a dominant inference cost. This motivates sharing key/value projections across heads.

In Multi-Query Attention (MQA), all heads share the same key and value projections, reducing KV cache size significantly. Grouped-Query Attention (GQA) is a compromise: queries remain fully multi-head, but keys/values are shared across groups of heads. In practice, GQA is common because it preserves some diversity while still yielding large memory savings.

\section{Transformer model families}
With these architectural components in mind, we can understand three major Transformer families that emerged from the original encoder--decoder design.

\subsection{Encoder--decoder models: T5 and span corruption}
Encoder--decoder Transformers are natural for sequence-to-sequence tasks such as translation. T5 (Text-to-Text Transfer Transformer) pushed a unifying idea: treat many NLP problems as text-in/text-out, and pretrain with an objective that differs from plain next-token prediction.

T5 uses span corruption. Instead of predicting the next token of a fully visible input, spans in the input are replaced with sentinel tokens. The encoder reads the corrupted input; the decoder generates the missing spans in sequence, using sentinel tokens to mark boundaries. This objective teaches the model to recover missing content while leveraging full context.

The T5 family includes multilingual variants (mT5) and byte-level variants (ByT5). ByT5 operates at the byte level with a small vocabulary (conceptually \(2^8\)), avoiding the typical large subword vocabulary. The trade-off is longer sequences and different inductive biases about how text is represented.

\subsection{Encoder-only models: BERT and bidirectional representations}
Encoder-only Transformers discard the decoder and are trained to produce strong contextual embeddings. BERT (Bidirectional Encoder Representations from Transformers) is the most influential example. “Bidirectional” here means each token can attend to both left and right context, unlike causal decoder attention.

BERT introduces special tokens that become conventions. A \texttt{[CLS]} token is placed at the beginning of the sequence; its final embedding is used as a sequence-level representation for classification tasks. A \texttt{[SEP]} token separates segments and supports paired-input tasks.

BERT pretraining combines two objectives. Masked Language Modeling (MLM) masks a subset of tokens and trains the model to predict them using bidirectional context. A practical masking recipe is to replace masked positions with \texttt{[MASK]} most of the time, leave some unchanged, and replace some with random tokens; this prevents the model from overfitting to the \texttt{[MASK]} symbol. Next Sentence Prediction (NSP) is a classification objective that predicts whether two segments are consecutive in the corpus. NSP relies on the \texttt{[CLS]} representation and was originally motivated as a way to teach cross-sentence coherence.

BERT’s input representation is the sum of token embeddings, positional embeddings, and segment embeddings. Segment embeddings are learned vectors indicating whether a token belongs to “sentence A” or “sentence B,” and all tokens within a segment share the same segment embedding. This is a simple but useful mechanism to mark structure for paired inputs.

After pretraining, BERT is fine-tuned by attaching a task-specific head. For sentiment classification, a classifier operates on the final \texttt{[CLS]} embedding. For token-level tasks such as question answering, heads operate across token embeddings to predict start and end spans.

\subsection{BERT variants: efficiency and recipe improvements}
Two influential refinements illustrate how training design choices matter as much as architecture.

DistilBERT compresses BERT via knowledge distillation. A large “teacher” model produces soft target distributions; a smaller “student” model is trained to match them (often using KL divergence). DistilBERT shows that reducing depth can yield substantial speedups while preserving much of the teacher’s performance, especially when distillation transfers the teacher’s “dark knowledge” in its probability distribution.

RoBERTa revisits BERT’s pretraining recipe. It shows that NSP can be removed without hurting—and sometimes improving—performance, and that dynamic masking (changing which tokens are masked across epochs) is beneficial. It also emphasizes that training longer and on more data can close gaps attributed to “undertraining” rather than architectural limitations.

\subsection{Decoder-only models: why they dominate LLMs}
Decoder-only models discard the encoder and rely on masked self-attention to generate text autoregressively. They align perfectly with next-token prediction on large corpora, which scales simply, and they map naturally onto chat and assistant use cases. While encoder-only models remain excellent for representation learning and classification, and encoder--decoder models remain strong for structured sequence-to-sequence tasks, decoder-only models became the default for general-purpose generation at scale.

\section{Summary}
Transformer design has evolved most dramatically in how it encodes position, how it normalizes activations, and how it manages attention cost at long context lengths. Relative and rotary position methods reflect the insight that attention depends primarily on relative distances. Pre-norm and RMSNorm reflect the need for stable deep optimization. Sparse attention and KV head sharing reflect the reality that memory and compute limits shape practical architectures.

At the model-family level, encoder--decoder systems such as T5 emphasize text-to-text learning via span corruption; encoder-only systems such as BERT emphasize bidirectional representations and classification-friendly embeddings; decoder-only systems dominate modern LLMs because they scale naturally with next-token prediction and support flexible generation.

\section{Exercises}
\paragraph{1. Absolute vs relative position.}
Explain why absolute positional embeddings can fail to extrapolate to longer contexts. Describe one advantage of RoPE or relative position bias in this setting.

\paragraph{2. Pre-norm stability.}
Give an intuitive argument for why applying normalization before a sublayer can improve stability in deep Transformers.

\paragraph{3. Attention complexity.}
If full attention is \(O(n^2)\) in sequence length \(n\), what is the approximate complexity of sliding-window attention with window size \(w\)? How can depth increase the effective receptive field?

\paragraph{4. GQA vs MQA.}
Why does sharing key/value projections reduce inference-time memory? What might you lose by making key/value projections too shared?

\paragraph{5. BERT objectives.}
Why does MLM require masking strategies beyond always replacing with \texttt{[MASK]}? What problem does NSP attempt to solve, and why might removing NSP still work in practice?
